% GENERATED FILE. Edit canonical lessons, then run npm run generate:books.
% canonical-source-hash: fnv1a64:d29c1d0cb7a02b10
% canonical-lessons: ML-C531-parihasikkuka, ML-C531-vimpilakkuka, ML-C531-anukarikkuka, ML-C531-ponnaccam, ML-C531-salyappetuttuka

\chapter{Mocking and Encouraging}
\label{ch:ml-mocking-and-encouraging}
\hlchaptermodality{531}

\begin{chapteropening}
\textbf{By the end of this chapter:} \emph{I can say to mock, to boast, to imitate, showing off and to bother.}

Complete the last lesson of chapter 531: I can say to mock, to boast, to imitate, showing off and to bother.
\end{chapteropening}

\section[\texorpdfstring{\ml{പരിഹസിക്കുക} (parihasikkuka)}{പരിഹസിക്കുക (parihasikkuka)}]{\ml{പരിഹസിക്കുക} (parihasikkuka) — to mock; to ridicule}

\label{lesson:ML-C531-parihasikkuka}



\begin{quote}
Before the new one: say the Malayalam for a coward, then the Malayalam for to feel shy.
\end{quote}

\subsection*{You'll want to know: \ml{പരിഹസിക്കുക}}
\textbf{\ml{പരിഹസിക്കുക}} — \emph{parihasikkuka} — \textquotedblleft{}to mock; to ridicule\textquotedblright{}.

\begin{quote}
\ml{അവന്റെ} \ml{ഉച്ചാരണത്തെ} \ml{കൂട്ടുകാർ} \ml{പരിഹസിച്ചു}.
\end{quote}

\begin{quote}
\emph{avanṟe uccāraṇatte kūṭṭukār parihasiccu.}
\end{quote}

\begin{quote}
His friends mocked his accent.
\end{quote}

Past: \ml{പരിഹസിച്ചു}, \emph{parihasiccu}.

Harsher than \ml{കളിയാക്കുക}; the noun is \ml{പരിഹാസം}, ridicule.

\begin{cousinweb}
\textbf{\ml{പരിഹസിക്കുക}} (\emph{parihasikkuka}) is built from Sanskrit \emph{parihāsa}, \textquotedblleft{}jesting, ridicule\textquotedblright{} (CDIAL 7902), from the root \emph{has}, \textquotedblleft{}to laugh\textquotedblright{}, with the Malayalam verb ending -\ml{ിക്കുക}.
\end{cousinweb}

\begin{grammarlens}[title={What you've built}]
One more word for describing what you lived and felt.
\end{grammarlens}

\subsection*{Your turn}
\begin{itemize}
\raggedright
  \item \emph{Say it:} \emph{parihasikkuka}
  \item \emph{Say it:} \emph{parihasikkuka}, once more
  \item \emph{Say it:} \emph{avanṟe uccāraṇatte kūṭṭukār parihasiccu.} — His friends mocked his accent.
  \item \emph{Recall:} say the Malayalam for a coward, then the Malayalam for to feel shy, then say \emph{parihasikkuka} again
\end{itemize}

\begin{tcolorbox}[breakable,colback=teal!4,colframe=teal!35!black,title={Before you move on}]
What does \ml{പരിഹസിക്കുക} mean? (\textquotedblleft{}To mock; to ridicule\textquotedblright{}.) Say it once more.
\end{tcolorbox}
\section[\texorpdfstring{\ml{വീമ്പിളക്കുക} (vīmpiḷakkuka)}{വീമ്പിളക്കുക (vīmpiḷakkuka)}]{\ml{വീമ്പിളക്കുക} (vīmpiḷakkuka) — to boast; to brag}

\label{lesson:ML-C531-vimpilakkuka}



\begin{quote}
Before the new one: say the Malayalam for to feel shy, then the Malayalam for to mock.
\end{quote}

\subsection*{You'll want to know: \ml{വീമ്പിളക്കുക}}
\textbf{\ml{വീമ്പിളക്കുക}} — \emph{vīmpiḷakkuka} — \textquotedblleft{}to boast; to brag\textquotedblright{}.

\begin{quote}
\ml{അവൻ} \ml{തന്റെ} \ml{പുതിയ} \ml{കാറിനെക്കുറിച്ച്} \ml{വീമ്പിളക്കി}.
\end{quote}

\begin{quote}
\emph{avan tanṟe putiya kāṟinekkuṟiccŭ vīmpiḷakki.}
\end{quote}

\begin{quote}
He boasted about his new car.
\end{quote}

Past: \ml{വീമ്പിളക്കി}, \emph{vīmpiḷakki}.

\ml{വീമ്പ്} a boast + \ml{ഇളക്കുക} to stir: to stir up one's own praise.

\begin{cousinweb}
\textbf{\ml{വീമ്പിളക്കുക}} (\emph{vīmpiḷakkuka}) is a Malayalam compound built from \ml{വീമ്പ്}, \textquotedblleft{}a boast\textquotedblright{}, whose origin is not settled (DEDR 5448 lists it, with a query, beside \ml{വീങ്ങുക}, \textquotedblleft{}to swell\textquotedblright{}), and \ml{ഇളക്കുക}, \textquotedblleft{}to stir, to shake\textquotedblright{}, which is native (DEDR 509).
\end{cousinweb}

\begin{grammarlens}[title={What you've built}]
One more word for describing what you lived and felt.
\end{grammarlens}

\subsection*{Your turn}
\begin{itemize}
\raggedright
  \item \emph{Say it:} \emph{vīmpiḷakkuka}
  \item \emph{Say it:} \emph{vīmpiḷakkuka}, once more
  \item \emph{Say it:} \emph{avan tanṟe putiya kāṟinekkuṟiccŭ vīmpiḷakki.} — He boasted about his new car.
  \item \emph{Recall:} say the Malayalam for to feel shy, then the Malayalam for to mock, then say \emph{vīmpiḷakkuka} again
\end{itemize}

\begin{tcolorbox}[breakable,colback=teal!4,colframe=teal!35!black,title={Before you move on}]
What does \ml{വീമ്പിളക്കുക} mean? (\textquotedblleft{}To boast; to brag\textquotedblright{}.) Say it once more.
\end{tcolorbox}
\section[\texorpdfstring{\ml{അനുകരിക്കുക} (anukarikkuka)}{അനുകരിക്കുക (anukarikkuka)}]{\ml{അനുകരിക്കുക} (anukarikkuka) — to imitate; to mimic}

\label{lesson:ML-C531-anukarikkuka}



\begin{quote}
Before the new one: say the Malayalam for to mock, then the Malayalam for to boast.
\end{quote}

\subsection*{You'll want to know: \ml{അനുകരിക്കുക}}
\textbf{\ml{അനുകരിക്കുക}} — \emph{anukarikkuka} — \textquotedblleft{}to imitate; to mimic\textquotedblright{}.

\begin{quote}
\ml{കുട്ടി} \ml{അച്ഛന്റെ} \ml{നടത്തം} \ml{അനുകരിച്ചു}.
\end{quote}

\begin{quote}
\emph{kuṭṭi acchanṟe naṭattaṁ anukariccu.}
\end{quote}

\begin{quote}
The child mimicked his father's walk.
\end{quote}

Past: \ml{അനുകരിച്ചു}, \emph{anukariccu}.

The noun is \ml{അനുകരണം}, imitation; the one imitated takes -\ml{എ}.

\begin{cousinweb}
\textbf{\ml{അനുകരിക്കുക}} (\emph{anukarikkuka}) is built from Sanskrit \emph{anukaraṇa}, \textquotedblleft{}imitation, following an example\textquotedblright{} (Monier-Williams), with the Malayalam verb ending -\ml{ിക്കുക}.
\end{cousinweb}

\begin{grammarlens}[title={What you've built}]
One more word for describing what you lived and felt.
\end{grammarlens}

\subsection*{Your turn}
\begin{itemize}
\raggedright
  \item \emph{Say it:} \emph{anukarikkuka}
  \item \emph{Say it:} \emph{anukarikkuka}, once more
  \item \emph{Say it:} \emph{kuṭṭi acchanṟe naṭattaṁ anukariccu.} — The child mimicked his father's walk.
  \item \emph{Recall:} say the Malayalam for to mock, then the Malayalam for to boast, then say \emph{anukarikkuka} again
\end{itemize}

\begin{tcolorbox}[breakable,colback=teal!4,colframe=teal!35!black,title={Before you move on}]
What does \ml{അനുകരിക്കുക} mean? (\textquotedblleft{}To imitate; to mimic\textquotedblright{}.) Say it once more.
\end{tcolorbox}
\section[\texorpdfstring{\ml{പൊങ്ങച്ചം} (poṅṅaccaṁ)}{പൊങ്ങച്ചം (poṅṅaccaṁ)}]{\ml{പൊങ്ങച്ചം} (poṅṅaccaṁ) — showing off; vanity}

\label{lesson:ML-C531-ponnaccam}



\begin{quote}
Before the new one: say the Malayalam for to boast, then the Malayalam for to imitate.
\end{quote}

\subsection*{You'll want to know: \ml{പൊങ്ങച്ചം}}
\textbf{\ml{പൊങ്ങച്ചം}} — \emph{poṅṅaccaṁ} — \textquotedblleft{}showing off; vanity\textquotedblright{}.

\begin{quote}
\ml{അവന്റെ} \ml{പൊങ്ങച്ചം} \ml{ആർക്കും} \ml{ഇഷ്ടമല്ല}.
\end{quote}

\begin{quote}
\emph{avanṟe poṅṅaccaṁ ārkkuṁ iṣṭamalla.}
\end{quote}

\begin{quote}
Nobody likes his showing off.
\end{quote}

\ml{പൊങ്ങച്ചം} \ml{പറയുക} is to boast; \ml{പൊങ്ങച്ചക്കാരൻ}, a show-off.

\begin{cousinweb}
\textbf{\ml{പൊങ്ങച്ചം}} (\emph{poṅṅaccaṁ}) is native Dravidian: DEDR 4469 lists it beside \ml{പൊങ്ങുക}, \textquotedblleft{}to boil over, to rise\textquotedblright{}, as a display of oneself (DEDR 4469).
\end{cousinweb}

\begin{grammarlens}[title={What you've built}]
One more word for describing what you lived and felt.
\end{grammarlens}

\subsection*{Your turn}
\begin{itemize}
\raggedright
  \item \emph{Say it:} \emph{poṅṅaccaṁ}
  \item \emph{Say it:} \emph{poṅṅaccaṁ}, once more
  \item \emph{Say it:} \emph{avanṟe poṅṅaccaṁ ārkkuṁ iṣṭamalla.} — Nobody likes his showing off.
  \item \emph{Recall:} say the Malayalam for to boast, then the Malayalam for to imitate, then say \emph{poṅṅaccaṁ} again
\end{itemize}

\begin{tcolorbox}[breakable,colback=teal!4,colframe=teal!35!black,title={Before you move on}]
What does \ml{പൊങ്ങച്ചം} mean? (\textquotedblleft{}Showing off; vanity\textquotedblright{}.) Say it once more.
\end{tcolorbox}
\section[\texorpdfstring{\ml{ശല്യപ്പെടുത്തുക} (śalyappeṭuttuka)}{ശല്യപ്പെടുത്തുക (śalyappeṭuttuka)}]{\ml{ശല്യപ്പെടുത്തുക} (śalyappeṭuttuka) — to bother; to pester}

\label{lesson:ML-C531-salyappetuttuka}



\begin{quote}
Before the new one: say the Malayalam for to imitate, then the Malayalam for showing off.
\end{quote}

\subsection*{You'll want to know: \ml{ശല്യപ്പെടുത്തുക}}
\textbf{\ml{ശല്യപ്പെടുത്തുക}} — \emph{śalyappeṭuttuka} — \textquotedblleft{}to bother; to pester\textquotedblright{}.

\begin{quote}
\ml{പഠിക്കുമ്പോൾ} \ml{അനിയൻ} \ml{എന്നെ} \ml{ശല്യപ്പെടുത്തി}.
\end{quote}

\begin{quote}
\emph{paṭhikkumpōḷ aniyan enne śalyappeṭutti.}
\end{quote}

\begin{quote}
My little brother bothered me while I was studying.
\end{quote}

Past: \ml{ശല്യപ്പെടുത്തി}, \emph{śalyappeṭutti}.

\ml{ശല്യം}, a nuisance, + -\ml{പ്പെടുത്തുക}, which makes a verb of causing it.

\begin{cousinweb}
\textbf{\ml{ശല്യപ്പെടുത്തുക}} (\emph{śalyappeṭuttuka}) is built on \ml{ശല്യം}, from Sanskrit \emph{śalya}, \textquotedblleft{}a dart, anything lodged in the body\textquotedblright{} (CDIAL 12352), with the native helper -\ml{പ്പെടുത്തുക}: a nuisance is a thorn in one's side.
\end{cousinweb}

\begin{grammarlens}[title={What you've built}]
One more word for describing what you lived and felt.
\end{grammarlens}

\subsection*{Your turn}
\begin{itemize}
\raggedright
  \item \emph{Say it:} \emph{śalyappeṭuttuka}
  \item \emph{Say it:} \emph{śalyappeṭuttuka}, once more
  \item \emph{Say it:} \emph{paṭhikkumpōḷ aniyan enne śalyappeṭutti.} — My little brother bothered me while I was studying.
  \item \emph{Recall:} say the Malayalam for to imitate, then the Malayalam for showing off, then say \emph{śalyappeṭuttuka} again
\end{itemize}

\begin{tcolorbox}[breakable,colback=teal!4,colframe=teal!35!black,title={Before you move on}]
What does \ml{ശല്യപ്പെടുത്തുക} mean? (\textquotedblleft{}To bother; to pester\textquotedblright{}.) Say it once more.
\end{tcolorbox}
